\documentclass[fleqn]{ctexart}
\usepackage{graphicx}
\usepackage{xcolor}
% 数学支持（提供 \mathbb、\dfrac 等）
\usepackage{amsmath,amssymb}
% 页面与版式设置：更小的页边距与顶部空白
\usepackage[a4paper,top=15mm,bottom=20mm,left=12mm,right=12mm]{geometry}
% 自动换行与字偶距优化
\usepackage{microtype}
\microtypesetup{tracking=true,protrusion=true,final}
% 在需要时放宽断行以避免溢出（数值可按需调整）
\setlength{\emergencystretch}{2em}
% 列表控制，便于让(1)后直接跟内容
\usepackage{enumitem}
% verbatim 环境支持（用于直接粘贴源码与输出）
\usepackage{verbatim}
\usepackage{fancyvrb}
\usepackage{float}
\usepackage{fvextra}
\RecustomVerbatimEnvironment{verbatim}{Verbatim}{
	frame=single,
	framesep=2mm,
	breaklines=true,
	breakanywhere=true
}
% 使章节标题左对齐（同时保留 ctex 默认样式）
\ctexset{section={format+=\raggedright}}
%1.3倍行距
\renewcommand{\baselinestretch}{1.3}
% 数学左对齐缩进为0
\setlength{\mathindent}{0pt}
% 增大矩阵的行距
\renewcommand{\arraystretch}{1.3}
\pagestyle{empty}
\begin{document}
\section*{题目1}

文件5\,MB，均分5段每段1\,MB在磁盘上连续存放。固定块方案块大小4\,KB，每个指针8\,B；extent方案每个extent的size域4\,B。

总块数 $=\frac{5\times1024}{4}=1280$ 块，每段 $=256$ 块，每间接块可存 $\frac{4096}{8}=512$ 个指针。

\subsection*{固定大小Data Block}

inode含12个直接指针、2个间接指针、1个双重间接指针。1280块按层级分配：

\begin{table}[H]
\centering
\caption{块分配}\label{tab:alloc}
\begin{tabular}{c|c|c|c}
\hline
指针 & 范围 & 块数 & 间接块 \\
\hline
12直接指针       & 0--11    & 12   & 0 \\
间接指针\#1     & 12--523  & 512  & 1 \\
间接指针\#2     & 524--1035 & 512  & 1 \\
双重间接        & 1036--1279 & 244 & 1+1 \\
\hline
合计            &          & 1280 & 4 \\
\hline
\end{tabular}
\end{table}

双重间接用了一个一级块和一个二级块覆盖剩余244块。元数据由间接块和inode内指针组成：

\begin{table}[H]
\centering
\caption{元数据}
\begin{tabular}{l|c}
\hline
4个间接块 $\times$ 4\,KB & 16384\,B \\
15个指针 $\times$ 8\,B   & 120\,B \\
\hline
合计                     & 16504\,B $\approx$ 16.1\,KB \\
\hline
\end{tabular}
\end{table}

\subsection*{Extent}

5段连续，每段一个extent。每个extent含起始指针8\,B和长度4\,B共12\,B，全部在inode内。$\;5\times12=60$\,B。

\subsection*{对比}

\begin{table}[H]
\centering
\caption{两种方案开销对比}
\begin{tabular}{l|c|c}
\hline
& 固定块 & Extent \\
\hline
额外间接块 & 4块\,/\,16\,KB & 0 \\
inode内占用 & 120\,B & 60\,B \\
\hline
总计 & 16504\,B & 60\,B \\
\hline
\end{tabular}
\end{table}

Extent方案60\,B，固定块方案约16.1\,KB，差距约275倍。

\section*{题目2}

VSFS布局：superblock、inode bitmap、data bitmap、inodes、data blocks。inode含12直接+2间接+1双重间接指针，指针4\,B，块大小4\,KB。每文件/目录至少分1块。操作独立，初始无缓存，访问一次即缓存。修改文件时该文件及其所在目录的访问/修改时间均更新，再上层不受影响，所有修改立刻落盘。

每间接块可存 $\frac{4096}{4}=1024$ 个指针。

\subsection*{操作1：/foo/bar追加1\,MB，原文件100\,KB}

原100\,KB占 $100/4=25$ 块，其中12块由直接指针覆盖，13块经间接指针。追加1\,MB即256块，共281块。间接块需要 $13+256=269<1024$ 个指针，不需要新间接块。

路径遍历需6次读：superblock $\to$ / inode $\to$ / 数据块 $\to$ foo inode $\to$ foo 数据块 $\to$ bar inode。

追加过程：读bar的间接块和data bitmap各1次；写256个数据块；写回data bitmap、间接块、bar inode和foo inode。

\begin{table}[H]
\centering
\caption{操作1 I/O汇总}
\begin{tabular}{c|c|l}
\hline
类型 & 次数 & 明细 \\
\hline
读 & 8 & 遍历6 + bar间接块 + data bitmap \\
写 & 260 & 数据块256 + bitmap + 间接块 + bar inode + foo inode \\
\hline
总计 & 268 & \\
\hline
\end{tabular}
\end{table}

\subsection*{操作2：mv /foo/bar /foo/bar2}

同目录重命名，路径遍历同上：superblock $\to$ / inode $\to$ / 数据块 $\to$ foo inode $\to$ foo 数据块，5次读。另需读bar inode以更新时间戳。修改阶段：写foo数据块将目录项bar改为bar2；写foo inode和bar inode更新mtime。不涉及bitmap和数据块。

\begin{table}[H]
\centering
\caption{操作2 I/O汇总}
\begin{tabular}{c|c|l}
\hline
类型 & 次数 & 明细 \\
\hline
读 & 6 & 遍历5 + bar inode \\
写 & 3 & foo数据块 + foo inode + bar inode \\
\hline
总计 & 9 & \\
\hline
\end{tabular}
\end{table}

\subsection*{对比}

\begin{table}[H]
\centering
\caption{两个操作I/O对比}
\begin{tabular}{l|c|c|c}
\hline
操作 & 读 & 写 & 总计 \\
\hline
追加1\,MB & 8 & 260 & 268 \\
mv重命名 & 6 & 3 & 9 \\
\hline
\end{tabular}
\end{table}

追加写入的I/O主要由数据块写入主导，元数据部分仅12次；同目录重命名不涉及数据块读写，仅9次I/O完成。

\end{document}